\clearpage
\thispagestyle{plain}
\phantomsection
\addcontentsline{toc}{chapter}{ABSTRAK}

\begin{center}
{\fontsize{12}{14}\selectfont\bfseries ABSTRAK}
\end{center}

\vspace{0.3cm}

\noindent
\textbf{\MakeUppercase{\JudulPA}}

\vspace{0.3cm}

\noindent
\textbf{Nama\hspace{1.45cm}: \NamaMahasiswa}\\
\textbf{NRP\hspace{1.65cm}: \NRP}\\
\textbf{Pembimbing\hspace{0.2cm}: \PembimbingSatu}\\
\textbf{\hspace{2.65cm} \PembimbingDua}

\vspace{0.3cm}

% TODO: Tulis abstrak Bahasa Indonesia di sini (200-300 kata)
Penyakit kardiovaskular merupakan penyebab kematian nomor satu di dunia. Deteksi dini aritmia melalui monitoring elektrokardiogram (ECG) secara kontinu menjadi sangat penting untuk memungkinkan intervensi medis yang tepat waktu. Proyek akhir ini mengimplementasikan sistem monitoring detak jantung berbasis IoT \textit{edge computing} menggunakan mikrokontroler ESP32-S3 dengan pendekatan TinyML untuk deteksi aritmia secara \textit{on-device}. Sistem yang dirancang meliputi algoritma \textit{preprocessing} sinyal ECG yang ringan, model TinyML terkuantisasi (INT8) untuk inferensi di perangkat, serta pipeline \textit{edge-to-cloud} yang resilien menggunakan protokol MQTT dengan mekanisme \textit{local buffering}.

\vspace{0.5cm}
\noindent
\textbf{Kata kunci:} ECG, IoT, Edge Computing, TinyML, ESP32-S3, Aritmia, MQTT

% ============================================================

\clearpage
\thispagestyle{plain}
\phantomsection
\addcontentsline{toc}{chapter}{ABSTRACT}

\begin{center}
{\fontsize{12}{14}\selectfont\bfseries ABSTRACT}
\end{center}

\vspace{0.3cm}

\noindent
\textbf{\MakeUppercase{\JudulPA}}

\vspace{0.3cm}

\noindent
\textbf{Name\hspace{1.2cm}: \NamaMahasiswa}\\
\textbf{NRP\hspace{1.45cm}: \NRP}\\
\textbf{Advisor\hspace{0.7cm}: \PembimbingSatu}\\
\textbf{\hspace{2.40cm} \PembimbingDua}

\vspace{0.3cm}

% TODO: Write English abstract here (200-300 words)
Cardiovascular disease is the leading cause of death worldwide. Early detection of arrhythmia through continuous electrocardiogram (ECG) monitoring is crucial for timely medical intervention. This final project implements a heart rate monitoring system based on IoT edge computing using the ESP32-S3 microcontroller with a TinyML approach for on-device arrhythmia detection. The designed system includes a lightweight ECG signal preprocessing algorithm, a quantized TinyML model (INT8) for on-device inference, and a resilient edge-to-cloud pipeline using MQTT protocol with a local buffering mechanism.

\vspace{0.5cm}
\noindent
\textbf{Keywords:} ECG, IoT, Edge Computing, TinyML, ESP32-S3, Arrhythmia, MQTT

\clearpage
